\section{A polynomial in twelve positive unknowns}\label{sec:twelve}

The supplied coordinates are exactly
\[
 \boxed{J,F,\alpha,z,h,i,T,s,w,\eta,\zeta,t\in\Pos.}
\]
The four former coordinates $\gamma,\tau,f,y$ are replaced by one coordinate
$t$. The polynomials in \eqref{eq:prefix}--\eqref{eq:units} are retained.

\subsection{The divisibility adjustment and the separating weights}

Define
\begin{align}
 G&=Hic^2,&D&=1+\Delta G^2,& Q&=\Delta(D-1)=(\Delta G)^2,
 &d&=Q-1,\notag\\
 v&=cT,&b&=c+RD,&e&=Q(Dv^2+b^2)-1,&B_0&=2Qvb.
 \label{eq:aux12}
\end{align}
The four equations whose witnesses will be eliminated are
\begin{equation}
 (P+H\gamma)^2=S,\qquad \tau^2=A,\qquad f^2=D,
 \qquad dy^2=e-B_0f.\label{eq:fourroots}
\end{equation}
The last equality is equivalent, when $f^2=D$, to
\[
 Q(vf-b)^2-(Q-1)y^2=1.
\]
This is the original normalized system with $i_0=Hi$, since
$vf-b=c(Tf-1)-Rf^2$. The factor $H$ in $G$ is the new divisibility
adjustment. Its role is the exact identity
\begin{equation}
 Q\equiv0\pmod{H^2},\qquad \gcd(H,d)=1.\label{eq:coprime}
\end{equation}
Completeness of this restriction will be proved below.

Use the positive polynomial weights
\begin{align}
 M&=1+8THic^2+16REH^2i^2c^3,\notag\\
 N&=1+2HM.\label{eq:weights12}
\end{align}
The coordinate $t$ is intended to encode
\begin{equation}
 t=y+M\tau+N\gamma.\label{eq:trace}
\end{equation}
In particular $N\equiv1\pmod H$. Neither weight is a new unknown.

\subsection{An explicit formula using only polynomial operations}

Set $W=Ht+NP$ and $m=HM$. Define the scalar polynomials
\begin{align}
 a_0&=d(W^2+m^2A+N^2S)-H^2e,\notag\\
 a_1&=-2dWm,&a_2&=-2dWN,&a_3&=2dmN,\label{eq:normcoeff1}\\
 f_0&=a_0^2+Aa_1^2+S(a_2^2+Aa_3^2)-H^4B_0^2D,\notag\\
 f_1&=2(a_0a_1+Sa_2a_3),\notag\\
 f_2&=2(a_0a_2+Aa_1a_3),\notag\\
 f_3&=2(a_0a_3+a_1a_2).\label{eq:normcoeff2}
\end{align}
Finally put
\begin{align}
 U&=f_0^2+Af_1^2-S(f_2^2+Af_3^2),\notag\\
 V&=2(f_0f_1-Sf_2f_3),\notag\\
 \Norm&=U^2-AV^2.\label{eq:norm12}
\end{align}
The requested single equation is
\begin{equation}
 \boxed{\quad P_{12}=N_kN_t(1+\Norm)-1=0.\quad}\label{eq:final12}
\end{equation}
An equivalent existence condition is the particularly transparent system
\begin{equation}
 \boxed{\quad N_k=1,\qquad N_t=1,\qquad \Norm=0.\quad}
 \label{eq:system12}
\end{equation}
Every formula here is an ordinary polynomial with integer coefficients.
The formulas contain no variable exponent, division, congruence test,
root operation, or hidden quantified parameter. The symbols used in the
proof for eliminated roots do not occur as unknowns in the final equation.
The emitted shared circuit computes $-f_1$ and $-f_2$
to avoid unnecessary negation gates, and factors
$H^4B_0^2D$ as $4(H^2Q)^2(Dv^2)b^2$.

The construction is two successive quadratic norms. In the algebra with
formal symbols $r,u$ satisfying $r^2=A$ and $u^2=S$, the expression
\[
 \bigl[d(W-mr-Nu)^2-H^2e\bigr]^2-H^4B_0^2D
\]
reduces to $f_0+f_1r+f_2u+f_3ru$. Taking its norm first in $u$ and then in
$r$ gives \eqref{eq:norm12}. This explanation is not part of the evaluator;
the scalar equations above define the polynomial directly.

\subsection{Separation of all algebraic branches}

The first task is to show that an integer root of the eliminated equation
does not merely encode a collection of irrational or rational nonintegral
witnesses. We prove separation before any Pell root is reconstructed.

\begin{lemma}[Separating inequalities]\label{lem:separation}
Under the positive input assumptions and the weak bounds
\eqref{eq:weak}, every algebraic branch of \eqref{eq:fourroots} satisfies
\[
 |y|<M\sqrt A,\qquad S>A,\qquad
 \frac{N\sqrt S}{H}>M\sqrt A+|y|.
\]
The absolute value is the complex absolute value; the first and last
inequalities include nonreal auxiliary branches.
\end{lemma}
\begin{proof}
All of $G,D,Q,v,b$ are positive and $Q>1$. Since $G>1$,
\[
 \sqrt D<(a+2)G<2EG.
\]
Put $Z_0=v\sqrt D+b>1$. For either sign of $f=\pm\sqrt D$,
\[
 |y|^2
 \le\frac{QZ_0^2+1}{Q-1}
 \le2Z_0^2+1<4Z_0^2.
\]
Thus $|y|<2(v\sqrt D+b)$. The right side is bounded above by
\[
 4cTEG+2c+8RE^2G^2.
\]
On the other hand, $\sqrt A>L=EkY>Ec/2$. Division by this positive
quantity gives
\begin{align*}
 \frac{|y|}{\sqrt A}
 &<8TG+\frac4E+16RE\frac{G^2}{c}\\
 &=8THic^2+\frac4E+16REH^2i^2c^3<M,
\end{align*}
since $E\ge16^4>4$.

For the second claim, $c>kY$ and $\Delta-E^2>2EY$ imply
\[
 \Delta c^2>k^2Y^2\Delta
 >k^2(E^2Y^2+EY)=A-1.
\]
Here $2EY^3>EY$, as $Y\ge1$. Consequently $S>A$. Finally
\[
 \frac{N\sqrt S}{H}>2M\sqrt S>2M\sqrt A>M\sqrt A+|y|.
\]
\end{proof}

\begin{lemma}[Injective weighted encoding]\label{lem:injective}
The map $(\gamma,\tau,f,y)\mapsto y+M\tau+N\gamma$ is injective
on the finite algebraic solution set of \eqref{eq:fourroots}.
\end{lemma}
\begin{proof}
If two tuples have different $\gamma$, the absolute difference of their
$N\gamma$ terms is $2N\sqrt S/H$. By Lemma~\ref{lem:separation}, this
strictly exceeds $2M\sqrt A+|y_1|+|y_2|$, which bounds the remaining
possible cancellation. If $\gamma$ agrees but $\tau$ differs, the
$M\tau$ difference has magnitude $2M\sqrt A>|y_1|+|y_2|$.
If both agree and the encoded value agrees, then $y$ agrees. The equality
$dy^2=e-B_0f$, with $B_0>0$, forces $f$ to agree as well.
\end{proof}

\subsection{Rationality and the integrality repair}

Over an algebraic closure, with multiplicities, the exact factorization is
\begin{equation}
 \Norm=d^8H^{16}
 \prod_{(\gamma,\tau,f,y)}
 \bigl(t-y-M\tau-N\gamma\bigr),\label{eq:factorization}
\end{equation}
where the sixteen branches are those of \eqref{eq:fourroots}.
Indeed, for fixed $r=\tau$ and $u=P+H\gamma$, the inner polynomial first
factors over the two values of $f$ and then over the two values of $y$;
each of its four linear factors contributes the indicated power of $H$,
and its leading coefficient in $t$ is $d^2H^4$. Multiplying over the four
choices of $r,u$ gives \eqref{eq:factorization}.

Suppose $\Norm=0$ at integer inputs satisfying the weak bounds. As $d,H$
are nonzero, one of the algebraic tuples in \eqref{eq:factorization} has
the specified integer value $t$. Choose a finite normal extension of
$\mathbb Q$ containing all its coordinates. Every automorphism over
$\mathbb Q$ sends it to another allowed tuple and fixes its encoded value.
Lemma~\ref{lem:injective} forces the whole tuple to be fixed. The fixed
field theorem therefore makes all four coordinates rational. This argument
does not assume that the square roots generate independent extensions.

Rationality is only the first half of the argument. Since
\[
 \tau^2=A,\qquad f^2=D,\qquad (P+H\gamma)^2=S
\]
are monic square equations with integer right sides, $\tau,f,u=P+H\gamma$
are integers. Hence the reduced denominator of $\gamma$ divides $H$.
If $y=r_0/s_0$ in lowest terms, then $dy^2=e-B_0f\in\Z$ implies
$s_0^2\mid d$. Thus the denominators of $\gamma$ and $y$ are coprime by
\eqref{eq:coprime}. But
\[
 N\gamma+y=t-M\tau\in\Z,\qquad \gcd(N,H)=1.
\]
Multiplication by $N$ cannot cancel a denominator dividing $H$, and the
remaining two coprime denominators cannot cancel each other. Both therefore
equal one. All four reconstructed witnesses are integers.

This last step is the reason for strengthening the Pell multiplier by $H$.
A bare resultant, or a rationality argument without this denominator
separation, would leave a real soundness gap.

\subsection{Soundness of the system and the single equation}

For \eqref{eq:system12}, the weak bootstrap holds directly. The preceding
argument restores \eqref{eq:fourroots}, and
\[
 vf-b=c(Tf-1)-Rf^2.
\]
The ordinary strong multiplier
\[
 i_{\rm ord}=\Delta Hi>0
\]
satisfies
\[
 (i_{\rm ord}c^2)^2=\Delta(f^2-1)=Q.
\]
The reconstructed tuple therefore satisfies the signed ordinary-strong
interface exactly. Its positive ports remain positive; $f,\tau,y,\gamma$
may be signed, as permitted by the direct signed theorem in
\src{upstream/signed_domain_core.md}. The witnesses $\tau,y$ cannot be zero:
$A>1$, while $y=0$ would give $Q(vf-b)^2=1$ with $Q>1$.
The inherited rank, extraction, and compiler proof gives a blank K215 halt.

For \eqref{eq:final12}, all three integer factors are units. By
\eqref{eq:norm12}, $V$ is even and
\[
 \Norm\equiv U^2\pmod4.
\]
Thus $1+\Norm$ cannot equal $-1$, whose residue is 3. It equals 1, so
$\Norm=0$ and $N_k=N_t=\varepsilon$ for a common sign. The weak bootstrap
was explicitly proved for these unit signs; the reconstruction consequently
remains valid. Applying the signed parent proof resolves
$\varepsilon=1$ and yields the same blank halt. No sign assumption was
silently used in the root elimination.

\subsection{Completeness of the strengthened multiplier}\label{sec:complete12}

Start from a canonical accepting blank tableau and its original first/main
Pell witnesses. Its actual rank satisfies $R\equiv3\pmod4$, and
$c=\psi_{a+2}(R)$. Put $A_0=a+2$. Consider the matrix
\[
 \mathcal M=\begin{pmatrix}A_0&\Delta\\1&A_0\end{pmatrix}.
\]
Its determinant is $A_0^2-\Delta=1$. Therefore its powers in the finite
group of invertible matrices modulo $Hc^2$ are periodic. Choose $m_0>0$
to be a sufficiently large common multiple of this period and $R$.
Then $\psi_{A_0}(m_0)$ is a positive multiple of $Hc^2$. Set
\begin{equation}
 i=\frac{\psi_{A_0}(m_0)}{Hc^2},\qquad
 f=\chi_{A_0}(m_0),\qquad
 \rho=\Delta\psi_{A_0}(m_0).\label{eq:periodic}
\end{equation}
All are positive integers and $\rho^2=Q$.

To construct the auxiliary witness explicitly, use
\begin{equation}
 V_* =\frac{\chi_\rho(R)}{\rho},\qquad y=\psi_\rho(R).
 \label{eq:auxcomplete}
\end{equation}
Because $R$ is odd, $\chi_x(R)$ is divisible by $x$ as an integer polynomial.
Write $R=2j+1$ and $\chi_x(R)=xQ_j(x^2)$. The inherited elementary
Chebyshev identities are
\[
 Q_j(0)=(-1)^jR,\qquad
 Q_j(1-A_0^2)=(-1)^j\psi_{A_0}(R).
\]
Now $j$ is odd. Moreover $\rho^2\equiv0\pmod c$ and
$\rho^2\equiv1-A_0^2\pmod f$. Consequently
\begin{equation}
 V_*\equiv-R\pmod c,\qquad V_*\equiv-c\pmod f.\label{eq:quotients}
\end{equation}
The quantities
\[
 o=\frac{V_*+c}{f},\qquad T=\frac{o+Rf}{c}
\]
are positive integers. To verify the second integrality claim, note that
$of=V_*+c\equiv-R\pmod c$, $f^2\equiv1\pmod c$, and $\gcd(f,c)=1$.
Hence $(o+Rf)f\equiv0\pmod c$, so $c\mid o+Rf$.
Substitution gives
\[
 c(Tf-1)-Rf^2=V_*.
\]
The Pell identity for \eqref{eq:auxcomplete} gives
$Q V_*^2-(Q-1)y^2=1$. All other canonical positive coordinates are
unchanged. Thus the stronger system has a positive witness whenever the
machine halts. No parity restriction on $m_0$ or on the new $i$ was needed.

Finally take $t=y+M\tau+N\gamma>0$ using the canonical positive roots.
One factor of \eqref{eq:factorization} vanishes, so both the system and
the single equation have a positive solution. Together with soundness and
the machine calibration, this proves Theorem~\ref{thm:main}.

\subsection{Exact degree and arithmetic size}

All supplied unknowns have degree one. The relevant degrees are
\begin{center}\small
\begin{tabular}{lrrrrrrrrrrrr}
\toprule
Expression &$c$&$E,a,H$&$\Delta$&$R$&$A,S$&$G$&$D$&$Q,d$&$b$&$M$&$N$\\
\midrule
Degree &5&6&12&4&22&17&46&58&50&39&45\\
\bottomrule
\end{tabular}
\end{center}
The exact total degree of $\Norm$ is 1360; its degree in $t$ is 16.
Here is a uniform certificate, rather than a purely syntactic degree bound.

Scale all twelve unknowns by a formal parameter, keeping the fixed
coefficients constant. The leading homogeneous forms satisfy
\[
 A_{\rm top}=S_{\rm top}=E_{\rm top}^{,2}c_{\rm top}^{,2}.
\]
The branches of $M\tau$ have leading terms $\pm16b_{\rm top}$ of degree 50.
The quadratic equation for $\gamma$ is
\[
 H^2\gamma^2+2HP\gamma+(P^2-S)=0,
\]
where the three coefficient degrees are 12, 17, and 16; the cancellation
\[
 P^2-S=X^2+2Xac-Hc^2-1
\]
is essential. Its branches have orders 5 and $-1$ in the scaling parameter.
The large branch contributes $-64b_{\rm top}$ to $N\gamma$ at degree 50;
the small branch contributes only degree 44. Each $y$ branch contributes
$\pm b_{\rm top}$ at degree 50, independently of the leading choice of $f$.
Thus the sixteen linear factors in \eqref{eq:factorization} have nonzero
leading magnitudes
\[
 15,17,15,17,15,17,15,17,
 47,49,79,81,47,49,79,81
\]
times $b_{\rm top}$. Equivalently,
\begin{equation}
 \Norm_{1360}=C_0d_{\rm top}^{,8}H_{\rm top}^{,16}b_{\rm top}^{,16},
 \qquad C_0=15^4 17^4(47\cdot49\cdot79\cdot81)^2>0.
 \label{eq:leader}
\end{equation}
The degree is $8\cdot58+16\cdot6+16\cdot50=1360$.
This may be interpreted as a formal Laurent-series calculation over the
field of rational functions in the unscaled coordinates. No numerical
approximation to an algebraic root is being used. All leading factors are
nonzero for every valid fixed coefficient binding with $\beta>0$.

The leading forms of $N_k,N_t$ have degrees 7 and 2 and are nonzero, so
$\deg P_{12}=1369$. The sum of the three squared residuals has exact degree
2720. The emitted live circuits have the following ledgers:
\begin{center}\small
\begin{tabular}{lrrrr}
\toprule
Presentation & Multiplications & Add/subtract & Total & Degree\\
\midrule
Three free comparisons &95&45&140&1360\\
Three explicit residuals &95&47&142&1360\\
Single product &97&47&144&1369\\
Sum of squares &98&49&147&2720\\
\bottomrule
\end{tabular}
\end{center}
The degree and coefficient in \eqref{eq:leader} are additionally checked
by a modular univariate specialization. The general leading-form argument
establishes uniformity on the actual astronomical coefficient slice.
